En un oscilador idealizado, el desplazamiento \(x(t)\), en metros, satisface
\[
x''+4x=0,
\]
con \(t\) en segundos; el coeficiente \(4\) tiene unidades \(\mathrm s^{-2}\). Se conoce \(x(0)=0.30\) m, pero se desconoce la velocidad inicial \(v=x'(0)\). Una lectura sintética exacta da \(x(\pi/6)=0.35\) m.

Recupere \(v\) a partir de esa lectura. Para estudiar el alcance de este procedimiento, exprese después una lectura general \(x(\tau)=z\), con \(\tau>0\), en función de \(v\). Determine en qué instantes permite recuperar una única velocidad y qué lecturas serían incompatibles, o compatibles con cualquier velocidad, en los demás instantes. Cuando la recuperación sea única, cuantifique cómo cambia el valor inferido de \(v\) al cambiar la lectura en una cantidad \(\eta\), y explique qué ocurre al acercar el instante de medida a uno que no distingue velocidades iniciales.
